\documentclass[11pt, a4paper]{article}

% ── Codificación y Lenguaje ──────────────────────────────────────────────────
\usepackage[utf8]{inputenc}
\usepackage[T1]{fontenc}
\usepackage[spanish, es-tabla]{babel}

% ── Geometría ───────────────────────────────────────────────────────────────
\usepackage[top=2.5cm, bottom=2.5cm, left=2.5cm, right=2.5cm, headheight=24pt]{geometry}

% ── Matemáticas ─────────────────────────────────────────────────────────────
\usepackage{amsmath, amssymb, amsthm}
\usepackage{mathtools}

% ── Colores y Cajas ─────────────────────────────────────────────────────────
\usepackage{xcolor}
\usepackage[most]{tcolorbox}
\tcbuselibrary{skins, breakable, theorems}

% ── Tablas ──────────────────────────────────────────────────────────────────
\usepackage{booktabs}
\usepackage{array}
\usepackage{multirow}
\usepackage{longtable}

% ── Listas ──────────────────────────────────────────────────────────────────
\usepackage{enumitem}

% ── Encabezado y Pie de Página ──────────────────────────────────────────────
\usepackage{fancyhdr}

% ── Secciones ───────────────────────────────────────────────────────────────
\usepackage{titlesec}

% ── Gráficos y Diagramas ────────────────────────────────────────────────────
\usepackage{tikz}
\usetikzlibrary{shapes.geometric, arrows.meta, positioning, calc}
\usepackage{graphicx}

% ── Columnas múltiples ──────────────────────────────────────────────────────
\usepackage{multicol}

% ── Espaciado ───────────────────────────────────────────────────────────────
\usepackage{setspace}
\setlength{\parskip}{4pt}
\setlength{\parindent}{0pt}

% ── Paleta GOP ──────────────────────────────────────────────────────────────
\definecolor{gopblue}{RGB}{31, 78, 121}
\definecolor{gopcyan}{RGB}{70, 130, 180}
\definecolor{gopgreen}{RGB}{0, 100, 0}
\definecolor{gopgreenlight}{RGB}{230, 255, 230}
\definecolor{goporange}{RGB}{200, 100, 0}
\definecolor{goporangelight}{RGB}{255, 245, 210}
\definecolor{gopred}{RGB}{160, 0, 0}
\definecolor{gopredlight}{RGB}{255, 230, 230}
\definecolor{gopgray}{RGB}{90, 90, 90}
\definecolor{gopgraylight}{RGB}{245, 245, 240}
\definecolor{gopbluedef}{RGB}{0, 70, 140}
\definecolor{gopbluedeflight}{RGB}{220, 235, 255}

% ── Hipervínculos y TOC ─────────────────────────────────────────────────────
\usepackage[colorlinks=true, linkcolor=gopblue, urlcolor=gopblue]{hyperref}

% ── Entornos tcolorbox ───────────────────────────────────────────────────────
\newtcolorbox{definicion}[1]{
  enhanced, breakable,
  colback=gopbluedeflight, colframe=gopbluedef,
  fonttitle=\bfseries\small, title={Definición: #1}, coltitle=white,
  attach boxed title to top left={yshift=-2mm, xshift=6pt},
  boxed title style={colback=gopbluedef, colframe=gopbluedef, sharp corners},
  sharp corners=south, top=4mm, before skip=6pt, after skip=6pt,
}
\newtcolorbox{formulas}[1]{
  enhanced, breakable,
  colback=gopgreenlight, colframe=gopgreen,
  fonttitle=\bfseries\small, title={Fórmulas: #1}, coltitle=white,
  attach boxed title to top left={yshift=-2mm, xshift=6pt},
  boxed title style={colback=gopgreen, colframe=gopgreen, sharp corners},
  sharp corners=south, top=4mm, before skip=6pt, after skip=6pt,
}
\newtcolorbox{tipprueba}[1]{
  enhanced, breakable,
  colback=goporangelight, colframe=goporange,
  fonttitle=\bfseries\small, title={TIP PRUEBA: #1}, coltitle=white,
  attach boxed title to top left={yshift=-2mm, xshift=6pt},
  boxed title style={colback=goporange, colframe=goporange, sharp corners},
  sharp corners=south, top=4mm, before skip=6pt, after skip=6pt,
}
\newtcolorbox{trampa}[1]{
  enhanced, breakable,
  colback=gopredlight, colframe=gopred,
  fonttitle=\bfseries\small, title={TRAMPA/ADVERTENCIA: #1}, coltitle=white,
  attach boxed title to top left={yshift=-2mm, xshift=6pt},
  boxed title style={colback=gopred, colframe=gopred, sharp corners},
  sharp corners=south, top=4mm, before skip=6pt, after skip=6pt,
}
\newtcolorbox{ejercicio}[1]{
  enhanced, breakable,
  colback=gopgraylight, colframe=gopgray,
  fonttitle=\bfseries\small\color{black}, title={Ejercicio: #1},
  attach boxed title to top left={yshift=-2mm, xshift=6pt},
  boxed title style={colback=gopgray!40, colframe=gopgray, sharp corners},
  sharp corners=south, top=4mm, before skip=6pt, after skip=6pt,
}

% ── Encabezado y pie ─────────────────────────────────────────────────────────
\pagestyle{fancy}
\fancyhf{}
\fancyhead[L]{\small\textbf{GOP ICS3213} --- Compendio de Ejercicios: Localización de Bodegas}
\fancyhead[C]{}
\fancyhead[R]{\small\thepage}
\fancyfoot[L]{\footnotesize Solución Oficial --- Transcripción en LaTeX}
\fancyfoot[R]{\small\thepage}
\renewcommand{\headrulewidth}{0.4pt}
\renewcommand{\footrulewidth}{0.4pt}

% ── Estilo de secciones ──────────────────────────────────────────────────────
\titleformat{\section}
  {\color{gopblue}\Large\bfseries}{\color{gopblue}\thesection.}{0.5em}{}
  [\color{gopblue}\titlerule]
\titleformat{\subsection}
  {\color{gopblue}\normalsize\bfseries}{\color{gopblue}\thesubsection.}{0.5em}{}
\titleformat{\subsubsection}
  {\color{gopblue}\small\bfseries}{\color{gopblue}\thesubsubsection.}{0.5em}{}

\begin{document}

% ════════════════════════════════════════════════════════════════════════════
% PORTADA
% ════════════════════════════════════════════════════════════════════════════
\begin{titlepage}
  \centering
  \vspace*{2cm}
  {\Huge\bfseries\color{gopcyan} Compendio de Ejercicios\par}
  \vspace{0.4cm}
  {\LARGE\bfseries Localización de Bodegas\par}
  \vspace{0.3cm}
  {\large Gestión de Operaciones (ICS3213)\par}
  \vspace{1.2cm}
  \rule{\textwidth}{0.4pt}
  \vspace{0.4cm}
  {\small Material extraído de pautas oficiales de Interrogaciones I2.\par}
  \vspace{0.4cm}
  \rule{\textwidth}{0.4pt}
  \vfill
\end{titlepage}

\newpage
\section*{Localización de Bodegas}

\subsection*{[I2_2020] Pregunta III: Localización de Bodegas y Punto de Equilibrio (20 Puntos)}
\begin{tipprueba}{¿Cuándo usar este enfoque?}
    \textbf{Identificación:} Se entregan coordenadas de fábricas (orígenes) y mercados (destinos) con demandas actuales y futuras (crecimiento lineal). Se pide calcular el Centro de Gravedad (CG), evaluar ubicaciones reales predefinidas y realizar un análisis de punto de equilibrio para elegir la mejor ubicación en un horizonte dinámico.\\
    \textbf{Qué aplicar:}
    \begin{itemize}
      \item Fórmulas del CG: $C_x = \frac{\sum V_i x_i}{\sum V_i}$, $C_y = \frac{\sum V_i y_i}{\sum V_i}$.
      \item Ecuación de Costo Total: $CT = CF + CV \cdot V$.
      \item Encontrar el punto de equilibrio en volumen: $CT_1(V) = CT_2(V)$.
      \item Evaluar el volumen dinámico lineal mediante áreas de costos acumulados sobre el horizonte temporal.
    \end{itemize}
  \end{tipprueba}

  \textbf{Enunciado:}
  Usted es dueño de una empresa que produce cemento. La empresa tiene dos fábricas (A y B) y distribuye a dos mercados (I y II). Actualmente debe decidir la ubicación de su centro de distribución y para ello ha recuperado la información de la producción de cada fábrica y consumo (en Toneladas de cemento) de cada mercado, para hoy y en 10 años más. Suponga un crecimiento lineal en las ventas.

  \begin{center}
  \small
  \begin{tabular}{lcccc}
    \toprule
    \textbf{Fábrica / Mercado} & \textbf{Prod/Cons HOY} & \textbf{Prod/Cons 10 Años} & \textbf{Coordenada X} & \textbf{Coordenada Y} \\
    \midrule
    \textbf{Fábrica A} & 20 & 260 & 10 & 10 \\
    \textbf{Fábrica B} & 70 & 140 & 40 & 20 \\
    \textbf{Mercado I} & 40 & 240 & 60 & 20 \\
    \textbf{Mercado II} & 50 & 160 & 70 & 80 \\
    \bottomrule
  \end{tabular}
  \end{center}

  Le han ofrecido a usted cuatro posibles ubicaciones para su centro de distribución, con sus respectivas coordenadas X e Y. Cada ubicación tiene un costo fijo de arriendo y un costo variable de transporte que depende únicamente de las toneladas totales de cemento que se muevan.

  \begin{center}
  \small
  \begin{tabular}{lcccc}
    \toprule
    \textbf{Ubicación} & \textbf{Coordenada X} & \textbf{Coordenada Y} & \textbf{Costo Fijo ($CF$)} & \textbf{Costo Variable ($CV$)} \\
    \midrule
    \textbf{Lugar 1} & 10 & 60 & 4000 & 10 \\
    \textbf{Lugar 2} & 42 & 29 & 6000 & 6 \\
    \textbf{Lugar 3} & 49 & 35 & 2000 & 16 \\
    \textbf{Lugar 4} & 58 & 20 & 1900 & 19 \\
    \bottomrule
  \end{tabular}
  \end{center}

  Con esta información:
  \begin{itemize}
    \item[\textbf{a)}] (6 Puntos) Determine la localización óptima de cada bodega para hoy y para 10 años usando el método del Centro de Gravedad.
    \item[\textbf{b)}] (6 Puntos) De las opciones de lugares propuestos determine el más adecuado para hoy y dentro de 10 años.
    \item[\textbf{c)}] (8 Puntos) Si debe elegir sólo una ubicación para los próximos 10 años. ¿Cuál elegiría? Muestre todos sus cálculos.
  \end{itemize}

  \medskip
  \begin{ejercicio}{Solución}
  \textbf{Desarrollo de la Solución:}

  \subsubsection*{a) Centro de Gravedad (CG)}
  Aplicamos las fórmulas del Centro de Gravedad:
  \begin{equation*}
    C_x = \frac{\sum V_i \cdot x_i}{\sum V_i}, \quad C_y = \frac{\sum V_i \cdot y_i}{\sum V_i}
  \end{equation*}

  \paragraph{Escenario HOY (Volumen Total = 180 Ton):}
  \begin{align*}
    C_x &= \frac{20 \cdot 10 + 70 \cdot 40 + 40 \cdot 60 + 50 \cdot 70}{180} = \frac{200 + 2800 + 2400 + 3500}{180} = \frac{8900}{180} \approx \textbf{49.44} \\
    C_y &= \frac{20 \cdot 10 + 70 \cdot 20 + 40 \cdot 20 + 50 \cdot 80}{180} = \frac{200 + 1400 + 800 + 4000}{180} = \frac{6400}{180} \approx \textbf{35.56}
  \end{align*}
  El Centro de Gravedad Hoy es $\mathbf{CG_{\text{Hoy}} = (49.44, 35.56)}$.

  \paragraph{Escenario 10 AÑOS (Volumen Total = 800 Ton):}
  \begin{align*}
    C_x &= \frac{260 \cdot 10 + 140 \cdot 40 + 240 \cdot 60 + 160 \cdot 70}{800} = \frac{2600 + 5600 + 14400 + 11200}{800} = \frac{33800}{800} \approx \textbf{42.25} \\
    C_y &= \frac{260 \cdot 10 + 140 \cdot 20 + 240 \cdot 20 + 160 \cdot 80}{800} = \frac{2600 + 2800 + 4800 + 12800}{800} = \frac{23000}{800} \approx \textbf{28.75}
  \end{align*}
  El Centro de Gravedad en 10 años es $\mathbf{CG_{\text{10Años}} = (42.25, 28.75)}$.

  \subsubsection*{b) Comparación de Opciones Propuestas}
  Comparamos las coordenadas teóricas con las ubicaciones reales evaluando su cercanía y costos totales:
  \begin{itemize}
    \item Para \textbf{HOY}, el $CG$ se ubica en $(49.44, 35.56)$. La opción física más cercana es el \textbf{Lugar 3} $(49, 35)$.
    \begin{equation*}
      \text{Costo Hoy (Lugar 3)} = CF_3 + CV_3 \cdot V_{\text{Hoy}} = 2000 + 16 \cdot 180 = \mathbf{\$4,880}
    \end{equation*}
    \item Para \textbf{10 AÑOS}, el $CG$ se ubica en $(42.25, 28.75)$. La opción física más cercana es el \textbf{Lugar 2} $(42, 29)$.
    \begin{equation*}
      \text{Costo 10 Años (Lugar 2)} = CF_2 + CV_2 \cdot V_{10} = 6000 + 6 \cdot 800 = \mathbf{\$10,800}
    \end{equation*}
  \end{itemize}

  \subsubsection*{c) Elección de Bodega Única a 10 Años (Punto de Equilibrio)}
  Buscamos el punto de corte en volumen ($V$) entre el Lugar 3 (costo fijo bajo) y el Lugar 2 (costo variable bajo):
  \begin{align*}
    CF_3 + CV_3 \cdot V &= CF_2 + CV_2 \cdot V \\
    2000 + 16 V &= 6000 + 6 V \\
    10 V &= 4000 \implies \mathbf{V = 400 \text{ Toneladas}}
  \end{align*}

  \begin{itemize}
    \item Si $V < 400$, es preferible el Lugar 3.
    \item Si $V > 400$, es preferible el Lugar 2.
  \end{itemize}

  El volumen crece linealmente desde 180 Ton (Año 0) hasta 800 Ton (Año 10):
  \begin{itemize}
    \item \textbf{Ahorro del Lugar 3 sobre el Lugar 2 (Volumen entre 180 y 400 Ton):}
    La variación de volumen es $\Delta V_1 = 400 - 180 = 220$ Ton.
    En $V=180$:
    \begin{align*}
      \text{Costo}_2(180) &= 6000 + 6 \cdot 180 = 7080 \\
      \text{Costo}_3(180) &= 2000 + 16 \cdot 180 = 4880 \\
      \text{Diferencia} &= 7080 - 4880 = 2200 \\
      \text{Área de Beneficio 3} &= \frac{220 \cdot 2200}{2} = \mathbf{\$242,000}
    \end{align*}

    \item \textbf{Ahorro del Lugar 2 sobre el Lugar 3 (Volumen entre 400 y 800 Ton):}
    La variación de volumen es $\Delta V_2 = 800 - 400 = 400$ Ton.
    En $V=800$:
    \begin{align*}
      \text{Costo}_3(800) &= 2000 + 16 \cdot 800 = 14800 \\
      \text{Costo}_2(800) &= 6000 + 6 \cdot 800 = 10800 \\
      \text{Diferencia} &= 14800 - 10800 = 4000 \\
      \text{Área de Beneficio 2} &= \frac{400 \cdot 4000}{2} = \mathbf{\$800,000}
    \end{align*}
  \end{itemize}

  \textbf{Conclusión:} Dado que el ahorro neto acumulado al estar en el \textbf{Lugar 2} ($\$800,000$) es significativamente superior al del \textbf{Lugar 3} ($\$242,000$), se debe seleccionar el \textbf{Lugar 2} como la ubicación definitiva para el horizonte de 10 años.

\end{ejercicio}

\newpage

\subsection*{[I2_2022] Opción A: Localización y Punto de Equilibrio Dinámico}
\textbf{Enunciado:}
  Usted debe decidir en qué lugar localizar su fábrica productiva y después de un análisis exhaustivo ha llegado a la conclusión de que sólo 2 lugares se ajustan a sus necesidades. Los dos lugares difieren en sus costos fijos y sus costos variables anuales, asociados al transporte por unidad. A continuación, se detallan los costos fijos y variables de cada ubicación:
  \begin{itemize}
    \item \textbf{LUGAR 1:} Costo Fijo = $\$10,000$, Costo Variable = $\$100/\text{unidad}$
    \item \textbf{LUGAR 2:} Costo Fijo = $\$23,000$, Costo Variable = $\$60/\text{unidad}$
  \end{itemize}

  Conteste:
  \begin{itemize}
    \item[\textbf{i.}] (5 ptos) Elabore un gráfico de punto de equilibrio y señale entre qué niveles productivos se hace conveniente cada lugar.
    \item[\textbf{ii.}] (7 ptos) Si en el año 1 la demanda es de 100 unidades y se espera que en 10 años llegue hasta 800 unidades, de forma lineal. ¿Cuál debería ser su decisión de ubicación? Muestre todos sus cálculos.
    \item[\textbf{iii.}] (8 ptos) Usted ha encargado un estudio de mercado, el cual le entrega un pronóstico de ventas. El estudio entrega un pronóstico de ventas para los primeros 5 años de $Q(T) = 100T$ para $T=1 \dots 5$ y de $Q(T) = 350 + 30T$ para $T=5 \dots 10$. Con esta información ¿cambia su decisión anterior? Muestre todos sus cálculos.
  \end{itemize}

  \medskip
  \begin{ejercicio}{Solución}
  \textbf{Desarrollo de la Solución:}

  \subsubsection*{i. Gráfico y Punto de Equilibrio}
  Buscamos el punto de corte de los costos totales:
  \begin{equation*}
    CF_1 + CV_1 \cdot Q = CF_2 + CV_2 \cdot Q
  \end{equation*}
  \begin{equation*}
    10000 + 100 Q = 23000 + 60 Q \implies 40 Q = 13000 \implies Q = 325 \text{ unidades}
  \end{equation*}

  \begin{itemize}
    \item Para $Q < 325$ unidades, el \textbf{Lugar 1} es óptimo (menores costos totales).
    \item Para $Q > 325$ unidades, el \textbf{Lugar 2} es óptimo.
  \end{itemize}

  \subsubsection*{ii. Crecimiento Lineal a 10 Años}
  La demanda crece de $Q_0 = 100$ (Año 1) a $Q_{10} = 800$ (Año 10).
  El punto de equilibrio se sitúa en $Q = 325$.
  Evaluamos la conveniencia integrando la diferencia de áreas geométricas de costo:
  \begin{itemize}
    \item \textbf{Rango de Conveniencia de Lugar 1 (de 100 a 325 unidades):}
      La variación de volumen es $\Delta Q_1 = 325 - 100 = 225$ unidades.
      En $Q = 100$:
      \begin{align*}
        \text{Costo}_1(100) &= 10000 + 100 \cdot 100 = 20000 \\
        \text{Costo}_2(100) &= 23000 + 60 \cdot 100 = 29000 \\
        \text{Diferencia} &= 29000 - 20000 = 9000 \\
        \text{Área de Ahorro Lugar 1} &= \frac{225 \cdot 9000}{2} = \$1,012,500
      \end{align*}

    \item \textbf{Rango de Conveniencia de Lugar 2 (de 325 a 800 unidades):}
      La variación de volumen es $\Delta Q_2 = 800 - 325 = 475$ unidades.
      En $Q = 800$:
      \begin{align*}
        \text{Costo}_1(800) &= 10000 + 100 \cdot 800 = 90000 \\
        \text{Costo}_2(800) &= 23000 + 60 \cdot 800 = 71000 \\
        \text{Diferencia} &= 90000 - 71000 = 19000 \\
        \text{Área de Ahorro Lugar 2} &= \frac{475 \cdot 19000}{2} = \$4,512,500
      \end{align*}
  \end{itemize}

  Dado que el beneficio neto esperado de elegir el \textbf{Lugar 2} ($\$4,512,500$) es mucho mayor que el del \textbf{Lugar 1} ($\$1,012,500$), la decisión correcta es elegir el \textbf{Lugar 2}.

  \subsubsection*{iii. Pronóstico de Ventas no Homogéneo}
  El pronóstico de demanda está fragmentado en dos tramos:
  \begin{itemize}
    \item \textbf{Tramo 1 (Años 1 a 5):} $Q(T) = 100T$
      \begin{itemize}
        \item Año 1: $Q(1) = 100$
        \item Año 5: $Q(5) = 500$
      \end{itemize}
    \item \textbf{Tramo 2 (Años 5 a 10):} $Q(T) = 350 + 30T$
      \begin{itemize}
        \item Año 5: $Q(5) = 500$
        \item Año 10: $Q(10) = 350 + 30 \cdot 10 = 650$
      \end{itemize}
  \end{itemize}

  Calculamos los costos anuales integrados para cada ubicación:

  \paragraph{Tramo 1 (Duración = 4 años, de T=1 a T=5)}
  \begin{itemize}
    \item \textbf{Lugar 1:}
      Costo en $Q=100$: $\$20,000$. Costo en $Q=500$: $\$60,000$.
      \begin{equation*}
        \text{Costo Total Tramo 1} = \frac{20000 + 60000}{2} \cdot 4 = \$160,000
      \end{equation*}
    \item \textbf{Lugar 2:}
      Costo en $Q=100$: $\$29,000$. Costo en $Q=500$: $\$53,000$.
      \begin{equation*}
        \text{Costo Total Tramo 1} = \frac{29000 + 53000}{2} \cdot 4 = \$164,000
      \end{equation*}
  \end{itemize}

  \paragraph{Tramo 2 (Duración = 5 años, de T=5 a T=10)}
  \begin{itemize}
    \item \textbf{Lugar 1:}
      Costo en $Q=500$: $\$60,000$. Costo en $Q=650$: $\$75,000$.
      \begin{equation*}
        \text{Costo Total Tramo 2} = \frac{60000 + 75000}{2} \cdot 5 = \$337,500
      \end{equation*}
    \item \textbf{Lugar 2:}
      Costo en $Q=500$: $\$53,000$. Costo en $Q=650$: $\$62,000$.
      \begin{equation*}
        \text{Costo Total Tramo 2} = \frac{53000 + 62000}{2} \cdot 5 = \$287,500
      \end{equation*}
  \end{itemize}

  \paragraph{Costos Totales Acumulados}
  \begin{itemize}
    \item \textbf{Lugar 1:} $\$160,000 + \$337,500 = \$497,500$
    \item \textbf{Lugar 2:} $\$164,000 + \$287,500 = \$451,500$
  \end{itemize}

  \textbf{Conclusión:} El \textbf{Lugar 2} sigue teniendo el menor costo total ($\$451,500 < \$497,500$), por lo que \textbf{no cambia} la decisión anterior.

\end{ejercicio}

\newpage

\end{document}
